% language: swedish
% figure: fig1 0.055 0.085 0.38 0.15
% figure: fig2 0.20 0.225 0.355 0.295
% figure: fig3 0.80 0.44 0.955 0.53
% figure: fig4 0.05 0.78 0.28 0.90
% figure: fig5 0.67 0.745 0.985 0.852
% figure: fig6 0.44 0.872 0.59 0.935
\section{Komplex Analys}
\hfill \rattat{Föreläsning 3/9}{Föreläsning 3/8}

\notefigure{fig1}{}
$f$ \emph{holomorf} om $\displaystyle\lim_{h \to 0} \frac{f(z+h) - f(z)}{h} = f'(z) \quad \forall z$

Ex: $z^3 - iz + 3,\ e^z,\ \sin z,\ \sum_{n=1}^{\infty} \frac{1}{n^z}, \dots$ \qquad ej $x, y$

\textbf{Integraler:}
\notefigure{fig2}{}
\begin{itemize}
  \item om $f$ holomorf: $\int_\gamma f(z)\,dz = 0$
  \item om $f$ singulär i enstaka punkter $\left(\frac{1}{1+z^2}\right)$ kan $\int_\gamma f(z)\,dz$ beräknas mha \emph{residykalkyl}
\end{itemize}
Låter oss även beräkna reella integraler, $\displaystyle\int_{-\infty}^{\infty} \frac{\cos x}{1+x^2}\,dx$, som i sin tur dyker upp i \emph{Fourieranalys}.

\textbf{Tillämpningar:} röntgenkristallografi (DNA), kvantmekanik, \dots

\textbf{Riemannhypotesen:} $\zeta(z) = \sum_{n=1}^{\infty} \frac{1}{n^z}$ kan vara noll på linjen
\notefigure{fig3}{}
bygger på \emph{identitetsprincipen}.

\noindent\rule{\linewidth}{0.4pt}

\subsection{Komplexa tal}
$\mathbb{C} = \{x + iy : x, y \in \mathbb{R}\}$, \quad $i$ nytt tal: $i^2 = -1$
\begin{align*}
  (x+iy) + (a+ib) &= (x+a) + i(y+b) \in \mathbb{C} \\
  (x+iy)(a+ib) &= (xa - yb) + i(xb + ya) \in \mathbb{C}
\end{align*}
Vanliga räkneregler gäller (Prop.\ 1)

\subsection{Komplexa talplanet}
\notefigure{fig4}{}
$x = \operatorname{Re} z$, \quad $y = \operatorname{Im} z$

$|z| = \sqrt{x^2 + y^2}$ \quad (absolutbelopp)

\notefigure{fig5}{}

\begin{framed}[Polär form:]
\notefigure{fig6}{}
$z = r(\cos\theta + i\sin\theta)$ \quad $\rightarrow$

$\arg z = \theta$ endast bestämt upp till multipel av $2\pi$
\end{framed}
